% Plantilla — el agente completa y guarda en presupuestos/<fecha-slug>/
% Compilar con XeLaTeX / LuaLaTeX / Tectonic. Ajustar \input@path según cwd.
\documentclass[11pt,a4paper]{article}
\makeatletter
\def\input@path{{../../latex/}{../latex/}{latex/}{./}}
\makeatother
\usepackage{tecnicocobre}
\hypersetup{
  pdftitle={Presupuesto — TITULO},
  pdfauthor={presupuesto-ar}
}
\begin{document}

\PresupuestoEncabezado{TITULO}{%
  Fecha de consulta: FECHA.\\
  Ámbito: precios de consumidor final en Argentina.}

\NotaIva

\section*{Detalle}
\begin{tabularx}{\linewidth}{@{}C L R R c@{}}
\cabeceracobre
\textbf{\#} & \textbf{Componente} & \textbf{Unitario} & \textbf{Subtotal} & \textbf{Ref.} \\
\midrule
1 & COMPONENTE & \$0 & \$0 & F01 \\
\bottomrule
\end{tabularx}

% \NoEncontrado{nombre exacto}

\TotalPresupuesto{\$0}

\section*{Fuentes}
\begin{itemize}
  \item \textbf{F01} · Mercado Libre · consultado FECHA · \url{https://www.mercadolibre.com.ar/}
\end{itemize}

{\footnotesize\color{muted}
Cada ítem se buscó por nombre exacto y se revalidó el precio en la ficha (2--3 lecturas). Los importes incluyen IVA según la publicación. Este PDF es referencia de mercado, no cotización formal ni orden de compra.
\par}

\end{document}
